\chapter{Analyse critique de l'état de l'art}

Les travaux présentés dans les premières sections de l’état de l’art mettent en évidence les difficultés 
particulières rencontrées par le traitement automatique des langues lorsqu’il s’applique à des langues peu 
dotées. Ces langues se caractérisent principalement par une disponibilité limitée de ressources numériques, 
notamment en ce qui concerne les corpus annotés, les lexiques bilingues ou les données parallèles nécessaires 
à l’entraînement des systèmes de traduction automatique. Cette situation constitue un obstacle majeur pour 
le développement d’outils de traitement linguistique performants, car les méthodes modernes reposent 
largement sur de grandes quantités de données d’apprentissage. Les langues régionales ou dialectales 
présentent en outre des défis supplémentaires liés à leur forte variabilité orthographique, à la coexistence 
de plusieurs variétés dialectales et à la dispersion des ressources disponibles. Malgré ces contraintes, 
l’étude de ces langues représente un domaine de recherche particulièrement important, dans la mesure où les 
méthodes développées pour ces contextes peuvent ensuite être adaptées à d’autres langues disposant de peu de 
ressources.

Dans ce contexte, la littérature distingue plusieurs types de corpus multilingues selon leur niveau 
d’alignement. Les corpus parallèles correspondent à des textes traduits dans différentes langues et alignés 
généralement au niveau de la phrase. Ils constituent historiquement la principale ressource utilisée dans 
les systèmes de traduction automatique statistique puis neuronale. Cependant, leur disponibilité reste 
limitée à certains domaines spécifiques, comme les textes institutionnels, juridiques ou administratifs. 
À l’inverse, les corpus comparables regroupent des documents traitant de sujets similaires dans différentes 
langues sans pour autant être des traductions directes les uns des autres. Ces corpus peuvent par exemple 
être constitués à partir d’articles encyclopédiques, d’articles de presse ou de contenus issus du web. 
Leur principal avantage réside dans leur relative facilité de collecte, notamment grâce aux ressources en 
ligne qui permettent d’obtenir de grandes quantités de textes dans plusieurs langues. L’utilisation de ces 
corpus repose en grande partie sur l’hypothèse distributionnelle selon laquelle des unités linguistiques 
apparaissant dans des contextes similaires dans différentes langues ont de fortes chances d’entretenir une 
relation de traduction ou de correspondance sémantique. Cette propriété permet d’exploiter indirectement des 
similarités entre langues même lorsque les textes ne sont pas explicitement alignés.

Toutefois, l’utilisation de corpus comparables introduit également plusieurs difficultés méthodologiques. 
L’absence d’alignement explicite entre les documents implique qu’il est nécessaire de recourir à des 
méthodes d’alignement automatique afin d’identifier les segments susceptibles d’être parallèles. De plus, 
ces corpus contiennent généralement un niveau de bruit plus élevé que les corpus parallèles, ce qui rend 
l’identification des correspondances plus complexe. Ces limitations expliquent l’importance des travaux 
consacrés à l’alignement de documents et de phrases dans des corpus multilingues.

Plusieurs approches ont été proposées dans la littérature pour répondre à ce problème. Certaines méthodes 
reposent sur l’extraction automatique de bitextes à partir du web. Des projets tels que ParaCrawl visent 
par exemple à explorer de grandes quantités de pages web afin d’identifier des documents traduits dans 
différentes langues. Ce processus implique généralement plusieurs étapes successives, comprenant la 
collecte de pages web, l’identification de documents potentiellement correspondants, l’alignement des
phrases et enfin le filtrage des données extraites afin d’éliminer les alignements incorrects. Des outils 
comme Hunalign utilisent pour cela des heuristiques relativement simples, fondées notamment sur la longueur 
des phrases ou la fréquence des mots, afin de déterminer la probabilité que deux segments correspondent à 
une traduction. Ces méthodes ont permis de constituer de grands corpus parallèles utilisés aujourd’hui pour  
l’entraînement de nombreux systèmes de traduction automatique. Néanmoins, elles supposent généralement 
l’existence de documents qui sont effectivement des traductions les uns des autres, ce qui limite leur 
applicabilité dans des contextes où les textes sont seulement comparables et non parallèles.

Face à ces limites, des approches plus récentes reposent sur l’utilisation de représentations vectorielles 
multilingues permettant de mesurer la similarité sémantique entre phrases dans différentes langues. Des 
modèles d’embeddings de phrases, comme ceux issus de travaux basés sur des architectures neuronales, 
permettent de projeter des segments textuels dans un espace vectoriel partagé où des phrases de sens 
similaire apparaissent proches les unes des autres, même lorsqu’elles appartiennent à des langues 
différentes. Ces représentations peuvent ensuite être utilisées pour identifier automatiquement des paires 
de phrases susceptibles d’être des traductions ou des paraphrases. Dans certains systèmes, ces scores de 
similarité sont combinés avec des méthodes d’alignement comme Vecalign, qui permettent de déterminer des 
correspondances entre segments de textes dans deux langues différentes. Les résultats obtenus peuvent 
ensuite être filtrés ou validés à l’aide de modèles de classification afin de distinguer les alignements 
corrects des correspondances erronées. Ces approches présentent l’avantage de ne pas dépendre exclusivement 
de la structure des documents ou de la longueur des phrases, mais de s’appuyer directement sur la 
proximité sémantique des segments.

Un autre axe de recherche concerne l’alignement documentaire, qui constitue souvent une étape préalable à 
l’alignement des phrases. Dans ce cadre, les méthodes inspirées du Cross-Lingual Information Retrieval 
visent à identifier des documents traitant du même sujet dans différentes langues. Le principe consiste 
à indexer un ensemble de documents dans une langue cible puis à utiliser des requêtes générées à partir 
de documents dans une autre langue afin de retrouver les textes les plus similaires. Des systèmes tels que 
Fast Document Aligner reposent sur ce type d’approche et permettent de sélectionner des documents candidats 
susceptibles de contenir des segments parallèles. Cette étape permet de réduire considérablement l’espace 
de recherche lors de l’alignement des phrases, en limitant l’analyse aux documents les plus pertinents. 
Toutefois, l’efficacité de ces méthodes dépend fortement de la qualité des ressources bilingues utilisées 
pour générer les requêtes ainsi que de la similarité thématique des documents comparés.

Dans l’ensemble, l’état de l’art montre que l’extraction de segments parallèles dans des corpus multilingues 
repose généralement sur une combinaison de plusieurs techniques complémentaires, incluant la constitution 
de corpus comparables, l’alignement documentaire et l’utilisation de modèles de similarité sémantique pour 
comparer les phrases. Ces approches ont principalement été développées pour des langues relativement bien 
dotées, ce qui signifie que leur application à des langues régionales ou dialectales reste encore 
relativement peu explorée. Dans ce contexte, le projet présenté dans ce mémoire s’inscrit dans une démarche 
visant à adapter ces méthodes aux spécificités des langues peu dotées. L’approche adoptée combine ainsi un 
alignement documentaire permettant d’identifier des textes traitant de sujets similaires avec 
l’utilisation de représentations vectorielles multilingues et d’un modèle de classification destiné à 
évaluer le degré de parallélisme entre segments. Cette combinaison de méthodes vise à tirer parti des 
avancées récentes du traitement automatique des langues tout en tenant compte des contraintes spécifiques 
liées à la disponibilité limitée des ressources linguistiques.